\begingroup
\small
\setlength{\tabcolsep}{3.5pt}
\renewcommand{\arraystretch}{1.15}
\begin{longtable}{@{}p{0.065\linewidth}p{0.285\linewidth}p{0.60\linewidth}@{}}
\caption{Impossible-tier task catalogue. Method set B applies to every task in this table.}\label{tab:catalogue-impossible}\\
\toprule
ID & Chemical task & Stated conditions and task-specific research target \\
\midrule
\endfirsthead
\multicolumn{3}{l}{\tablename\ \thetable\ (continued)}\\
\toprule
ID & Chemical task & Stated conditions and task-specific research target \\
\midrule
\endhead
\midrule
\multicolumn{3}{r}{Continued on the next page}\\
\endfoot
\bottomrule
\endlastfoot
IA1 & cyclobutadiene D4h automerization\newline {\footnotesize Charge 0; multiplicity 1.} & \textit{Conditions:} gas phase, low temperature (matrix isolation regime).\newline \textit{Target:} Thermal automerization of rectangular cyclobutadiene through the square transition structure: mechanism and activation barrier. \\
\addlinespace[3pt]
IA5 & cyclobutene disrotatory ring opening\newline {\footnotesize Charge 0; multiplicity 1.} & \textit{Conditions:} gas phase, thermal.\newline \textit{Target:} Thermal disrotatory ring opening of cyclobutene to butadiene: mechanism and activation barrier. \\
\addlinespace[3pt]
IA7 & Bergman cyclization of (Z)-hex-3-ene-1,5-diyne\newline {\footnotesize Charge 0; multiplicity 1.} & \textit{Conditions:} gas phase, thermal.\newline \textit{Target:} Bergman cyclization to p-benzyne: mechanism and activation barrier. \\
\addlinespace[3pt]
IA9 & 1,2-dioxetane thermolysis to two formaldehyde\newline {\footnotesize Charge 0; multiplicity 1.} & \textit{Conditions:} gas phase, thermal.\newline \textit{Target:} Thermolysis of 1,2-dioxetane to two formaldehyde molecules: mechanism and activation barrier. \\
\addlinespace[3pt]
IB9 & diazomethane and ethylene to cyclopropane via 1-pyrazoline\newline {\footnotesize Charge 0; multiplicity 1.} & \textit{Conditions:} gas-phase thermolysis.\newline \textit{Target:} Conversion of diazomethane and ethylene to cyclopropane and nitrogen through 1-pyrazoline. \\
\addlinespace[3pt]
IC3 & phenylnitrene open-shell singlet state\newline {\footnotesize Charge 0; multiplicity 3.} & \textit{Conditions:} gas phase.\newline \textit{Target:} Singlet--triplet energy gap and the electronic character of the lowest phenylnitrene singlet state. \\
\addlinespace[3pt]
ID2 & anomalous bonding in dioxygen difluoride (FOOF)\newline {\footnotesize Charge 0; multiplicity 1.} & \textit{Conditions:} gas phase, low temperature.\newline \textit{Target:} Equilibrium F--O and O--O distances of FOOF and its F--O bond-dissociation energy. \\
\addlinespace[3pt]
IE2 & nitrogen triple-bond dissociation curve\newline {\footnotesize Charge 0; multiplicity 1.} & \textit{Conditions:} gas phase.\newline \textit{Target:} Ground-state N2 dissociation curve and bond-dissociation energy. \\
\addlinespace[3pt]
IF3 & triangulene triplet ground state and singlet-triplet gap\newline {\footnotesize Charge 0; multiplicity 3.} & \textit{Conditions:} gas phase (surface-synthesis analogue).\newline \textit{Target:} Ground-state spin multiplicity and singlet--triplet gap of triangulene. \\
\addlinespace[3pt]
IG8 & chromium dimer ground-state dissociation curve\newline {\footnotesize Charge 0; multiplicity 1.} & \textit{Conditions:} gas phase.\newline \textit{Target:} Ground-state Cr2 potential-energy curve and bond-dissociation energy. \\
\addlinespace[3pt]
\end{longtable}
\endgroup
